\documentclass[12pt]{article}
\usepackage[utf8]{inputenc}
\usepackage[T1]{fontenc}
\usepackage{hyperref}
\title{Tatuyo Language Resources for K-12 Curriculum Development}
\author{Compiled by Claude Sonnet 5 (claude-sonnet-5) for Andr\'as Kornai}
\date{}
\begin{document}
\maketitle
\section{Introduction}
Tatuyo (ISO 639-3 tav; Glottocode tatu1247) is an East Tucanoan language of Colombia with about 400 speakers. Its lexicon is reported to be about 96\% cognate with Carapana, its closest relative. Glottolog lists it as endangered and shifting. The search found no published Tatuyo primer, textbook, dictionary or grammar.
\section{Primary Educational Resources}
None were found in search results. Grammars of the related but distinct language Tukano (tuo), for example West's pedagogical grammar, exist, but they are not Tatuyo materials.
\section{Linguistic and Reference Materials}
\begin{itemize}
\item Wikipedia, Tatuyo language (overview: six vowels, tonal, Carapana relationship): \url{https://en.wikipedia.org/wiki/Tatuyo_language}
\item Glottolog languoid page for Tatuyo: \url{https://glottolog.org/resource/languoid/id/tatu1247}
\item ASJP Tatuyo wordlist (comparative word list, not a dictionary): \url{https://asjp.clld.org/languages/TATUYO}
\item WALS entry for Tatuyo: \url{https://wals.info/languoid/lect/wals_code_tty}
\end{itemize}
\section{Religious and Community Resources}
No Tatuyo scripture or community organisation was found in the searches.
\section{Digital and Audio Resources}
No Tatuyo audio, video or software resources were found.
\section{Curriculum Development Context}
Very little documentation exists. A K-12 curriculum would require new fieldwork and orthography work with the Tatuyo community. Related Tucanoan materials, such as Tukano grammars cataloged at \url{https://glottolog.org/resource/reference/id/92705}, may offer comparative guidance only.
\end{document}